% Auto-generated by src/analysis/latex_tables.py — do not edit by hand.
\begin{table}[htbp]
\centering
\small
\caption{Magnitude decomposition: Kulturkampf-attributable change vs.\ the observed differential between high- and low-Catholic counties}
\label{tab:magnitudes}
\begin{tabular}{lccccc}
\toprule
 & \multicolumn{3}{c}{Observed change post $-$ pre} & \multicolumn{2}{c}{Decomposition} \\
\cmidrule(lr){2-4}\cmidrule(lr){5-6}
 & High-Cath & Low-Cath & Differential & IV-implied & Counterfactual \\
 & ($>$75\%) & ($<$25\%) & (1)$-$(2) & Kulturkampf & gap (3)$-$(4) \\
 & (1) & (2) & (3) & (4) & (5) \\
\midrule
Crude birth rate & -0.229 & -0.280 & +0.050 & -3.489 & +3.539 \\
Legit.\ birth rate & +0.004 & -0.144 & +0.149 & -4.219 & +4.368 \\
Illegitimacy ratio & -0.444 & -0.359 & -0.085 & +2.753 & -2.838 \\
Marriage rate & -0.618 & -0.355 & -0.263 & -0.665 & +0.402 \\
\bottomrule
\end{tabular}
\begin{minipage}{\linewidth}
\vspace{0.5em}
\footnotesize \textit{Notes:} Columns (1)--(2) report the change in the outcome between the 1862--71 mean and the 1880--89 mean for counties above 75\% Catholic and below 25\% Catholic respectively. Column~(3) is the differential change. Column~(4) is the Kulturkampf-attributable differential implied by the 2SLS coefficient: $\hat{\beta}_{IV} \times (\bar{\text{cath}}_{\text{high}} - \bar{\text{cath}}_{\text{low}})$. Column~(5) is the counterfactual differential absent the Kulturkampf, $(3) - (4)$. A negative Kulturkampf-attributable column with a positive observed differential (as for CBR) implies that secular pre-trends were widening the high--low gap and that the Kulturkampf partially offset, rather than drove, the observed pattern.
\end{minipage}
\end{table}
